\scanpage{84}
% loc 1–5
Finalement, on trouve une ``table de multiplication'' pour une $k$-algèbre libre de rang 4 (dont $V$ sera quotient) \unsure{explicite|explicitée} en termes de 5 éléments de $k$
% loc 6–8, display and notes at right
\[
\left\{\begin{aligned}
 b&=\operatorname{Tr}u,\ c=\det u\deleted{\unsure{,a}}&
 &\text{(permet de décrire } k[u]\inserted{=k\oplus k.u}\\
 &&&\quad \inserted{\text{(si }k\cap k.u=\{0\}\text{)}}\text{)}\\
 b'&=\operatorname{Tr}v,\ c'=\det v\deleted{\unsure{=}}&&\text{(id pour } k[v]\text{)}\\
 t&=\operatorname{Tr}(uv)
\end{aligned}\right. \tag*{(7)}
\]
% loc 9–12
avec comme base $\underline{1},\underline{u},\underline{v},\underline{uv}$ (les soulignés indiquent qu'on considère $1,u\ \unsure{..}$ comme éléments de la nouvelle algèbre). On appelle cette Algèbre
% loc 13
\[ M(\deleted{\unsure{a}}b,\deleted{\unsure{b}}c;\deleted{\unsure{a'}}b',\deleted{\unsure{b'}}c';t) \tag*{(8)} \]
% loc 14–16
On vérifie qu'elle existe (indépendamment de la donnée de $A$) pour tout choix de \deleted{$b,c,$}
% loc 17
\[ (b,c,b',c',t)\in k^5 \tag*{(\unsure{4|9})} \]
% loc 18–19
En fait, on a une alg. canonique \unsure{de} rang 4 sur
% loc 20
\[ \mathbb Z[B,C,B',C',T] \]
% loc 21
où $B,C,B',C',T$ sont des indéterminées.

% loc 22–24
\uline{Vérification} On écrit une table de multiplication explicite, en utilisant H C pour \unsure{$u,v$|$u$ et $v$} et (6) alias $(6')$ (pour $(u,v)$),
% archive sheet number: 83
